\ifdefined\mainfile\else\documentclass{article}% ============================================================
%  FÆLLES PRÆAMBEL
%  Deles af main.tex (hele rapporten) og de enkelte sektionsfiler,
%  som via \ifdefined\mainfile også kan kompileres alene.
%  Ret kun pakker/stil her ét sted.
%  NB: \documentclass er IKKE her -- den står i main.tex (og i sektionernes
%  standalone-guard), så Overleaf ikke forveksler preamble.tex med hoveddokumentet.
% ============================================================
\usepackage[utf8]{inputenc}
\usepackage[T1]{fontenc}
\usepackage[danish]{babel}
\usepackage{subcaption}
\usepackage{graphicx}              % Indsættelse af billeder
\usepackage[absolute,overlay]{textpos}
\usepackage{float}                 % H-placering
\usepackage{amsmath}
\usepackage{amssymb}
\usepackage{siunitx}               % \SI{15}{\volt} m.m. -- praktisk til måleresultater
\usepackage{booktabs}
\usepackage[a4paper, left=2.5cm, right=2.5cm, top=2.5cm, bottom=2.5cm]{geometry}
\usepackage{listings}              % Kildekode-listings
\usepackage{xcolor}
\usepackage{tikz}
\usetikzlibrary{automata, positioning, arrows, calc, shapes.geometric}
\tikzset{
    >=stealth,
    shorten >=2pt, shorten <=2pt,
    node distance=2.8cm,
    block/.style={draw, very thick, rectangle, minimum height=1cm, minimum width=2cm, fill=blue!10},
}
\usepackage[colorlinks=true,
            linkcolor=black,
            citecolor=blue,
            urlcolor=blue]{hyperref}

% ------------------------------------------------------------
%  Listing-stil til C / Arduino-firmware
% ------------------------------------------------------------
\lstdefinestyle{c}{
    language=C,
    basicstyle=\footnotesize\ttfamily,
    keywordstyle=\color{blue}\bfseries,
    commentstyle=\color{green!60!black}\itshape,
    stringstyle=\color{red},
    numbers=left, numberstyle=\tiny\color{gray}, stepnumber=1, numbersep=8pt,
    showspaces=false, showstringspaces=false, showtabs=false,
    frame=single, rulecolor=\color{black}, tabsize=4, captionpos=b,
    breaklines=true, breakatwhitespace=false,
    literate=%
        {æ}{{\ae}}1 {ø}{{\o}}1 {å}{{\aa}}1
        {Æ}{{\AE}}1 {Ø}{{\O}}1 {Å}{{\AA}}1,
}

% ------------------------------------------------------------
%  Listing-stil til MATLAB
% ------------------------------------------------------------
\lstdefinestyle{matlab}{
    language=Matlab,
    basicstyle=\footnotesize\ttfamily,
    keywordstyle=\color{blue}\bfseries,
    commentstyle=\color{green!60!black}\itshape,
    stringstyle=\color{magenta},
    numbers=left, numberstyle=\tiny\color{gray}, numbersep=8pt,
    showstringspaces=false, frame=single, tabsize=4, captionpos=b, breaklines=true,
}

\title{Elektrisk Energisystem}
\author{Gruppe \textbf{XX} -- 62768}
\date{Juni 2026}

% \mainfile defineres i main.tex (IKKE her) -- ellers tror standalone-sektioner
% også at de er i hoveddokumentet og springer deres \end{document} over.
\newcommand{\doubleunderline}[1]{\underline{\underline{#1}}}
\begin{document}\fi

\subsection{Current measurements}
\label{sec:stroemmaaling}
% Strømmåling -- op-amps tilladt, ingen øvrige IC'er (krav 10). INA219 anvendt efter aftale.

\noindent
Current measurement is needed in two places in the system. First, the MPPT algorithm
needs the panel current in order to compute the PV power $P = V \cdot I$ and track the
maximum power point (see Section~\ref{sec:styring}). Second, the supply currents are
used to verify the operation of the complete system, in particular to show that the
solar panel offloads the motor--generator.

\vspace{0.5cm} \noindent
The underlying measurement principle is the same in all cases. A small \emph{shunt
resistor} is placed in series with the current path, and the voltage drop across it is
measured. By Ohm's law the current then follows directly as
\[
I = \frac{V_{\text{shunt}}}{R_{\text{shunt}}}.
\]
Because the shunt resistance is small, the inserted voltage drop is kept low, so the
measurement does not noticeably disturb the circuit. The shunt resistor itself is a
single passive component, and the differential voltage across it is what has to be
amplified and digitised.

\vspace{0.5cm} \noindent
On the PV side this measurement is implemented with an INA219 module, which contains a
$0.1~\Omega$ shunt resistor together with an instrumentation amplifier and a 12-bit
analog-to-digital converter. The module is configured for a $32~\text{V}$ bus range and
a $\pm 320~\text{mV}$ shunt range (gain $\div 8$), and its calibration register is set so
that the current and power registers read directly in amps and watts. The Arduino reads
the bus voltage, the shunt voltage and the resulting current from the INA219 over the
I\textsuperscript{2}C bus (pins A4/A5). The series shunt resistor, whose differential
voltage is amplified and converted, is the core of the measurement, and the same shunt
could equally be read by an op-amp difference amplifier ahead of the microcontroller
ADC.

\vspace{0.5cm} \noindent
In the MPPT control loop the INA219 returns the panel voltage and current in every
iteration, and these are multiplied to obtain the instantaneous PV power used by the
Perturb-and-Observe algorithm. The measured panel currents span the full operating
range of the module, from approximately $0.60~\text{A}$ near short circuit down to about
$0.05~\text{A}$ close to the open-circuit voltage, as listed in
Table~\ref{tab:pv_measurements}. The supply currents for the complete-system test were
read directly from the laboratory power supplies: the motor drew $1.412~\text{A}$ with
the solar panel inactive and $0.944~\text{A}$ with it active, a reduction of about
$33\%$, which confirms that the panel contributes power and reduces the load on the
motor--generator.

\ifdefined\mainfile\else\end{document}\fi
